\begin{figure}[tbp]
\centering
\begin{tikzpicture}[
  every node/.style={font=\footnotesize},
  live/.style={fill=allocring!30},
  done/.style={pattern=north east lines, pattern color=allocring!70!black},
  gap/.style={pattern=crosshatch dots, pattern color=gray!70},
  cur/.style={-{Stealth[length=5pt]}, very thick},
  headc/.style={orange!85!black},
  tailc/.style={blue!70!black},
]
% One contiguous pool, drawn growing upwards; each column is a snapshot.
% Heights are in arbitrary units, 8 for the whole pool.
\newcommand{\rgU}{0.45}
\newcommand{\rgW}{1.1}
\newcommand{\rgpool}[1]{\draw[fill=gray!6] (#1,0) rectangle (#1+\rgW,8*\rgU);}
\newcommand{\rgblock}[5]{% x, from, to, style, label
  \draw[#4] (#1,#2*\rgU) rectangle (#1+\rgW,#3*\rgU);
  \node at (#1+\rgW/2,{(#2+#3)/2*\rgU}) {#5};}
\newcommand{\rghead}[2]{\draw[cur, headc] (#1+\rgW+0.45,#2*\rgU) -- (#1+\rgW+0.02,#2*\rgU);}
\newcommand{\rgtail}[2]{\draw[cur, tailc] (#1-0.45,#2*\rgU) -- (#1-0.02,#2*\rgU);}
\newcommand{\rgstep}[2]{\node[align=center, anchor=north, text width=2.3cm] at (#1+\rgW/2,-0.12) {#2};}
\newcommand{\rgB}{{\fboxsep=1pt\colorbox{white}{B}}}

% (a) empty
\rgpool{0}
\rgtail{0}{0}\rghead{0}{0}
\rgstep{0}{(a) empty}

% (b) three allocations of different sizes, back to back
\rgpool{2.4}
\rgblock{2.4}{0}{1.2}{live}{A}\rgblock{2.4}{1.2}{3.5}{live}{B}\rgblock{2.4}{3.5}{5.4}{live}{C}
\rgtail{2.4}{0}\rghead{2.4}{5.4}
\rgstep{2.4}{(b) allocate\\A, B, C}

% (c) free B: marked done, but A is older and still live
\rgpool{4.8}
\rgblock{4.8}{0}{1.2}{live}{A}\rgblock{4.8}{1.2}{3.5}{done}{\rgB}\rgblock{4.8}{3.5}{5.4}{live}{C}
\rgtail{4.8}{0}\rghead{4.8}{5.4}
\rgstep{4.8}{(c) free B\\{\scriptsize\itshape tail stays}}

% (d) free A: tail moves past A and B at once
\rgpool{7.2}
\rgblock{7.2}{3.5}{5.4}{live}{C}
\rgtail{7.2}{3.5}\rghead{7.2}{5.4}
\rgstep{7.2}{(d) free A\\{\scriptsize\itshape tail skips B too}}

% (e) allocate D (3.0): only 2.6 left above the head, so wrap
\rgpool{9.6}
\rgblock{9.6}{0}{3.0}{live}{D}\rgblock{9.6}{3.5}{5.4}{live}{C}\rgblock{9.6}{5.4}{8}{gap}{}
\rgtail{9.6}{3.5}\rghead{9.6}{3.0}
\rgstep{9.6}{(e) allocate D\\{\scriptsize\itshape head wraps}}

% (f) deallocate_all
\rgpool{12.0}
\rgtail{12.0}{0}\rghead{12.0}{0}
\rgstep{12.0}{(f) reset\\{\scriptsize\texttt{deallocate\_all}}}

% legend
\begin{scope}[shift={(-0.4,-1.55)}]
  \draw[cur, tailc] (0,0.125) -- (0.45,0.125); \node[right] at (0.5,0.125) {tail};
  \draw[cur, headc] (1.45,0.125) -- (1.9,0.125); \node[right] at (1.95,0.125) {head};
  \draw[live] (3.2,0) rectangle (3.6,0.25); \node[right] at (3.65,0.125) {live};
  \draw[done] (4.75,0) rectangle (5.15,0.25); \node[right] at (5.2,0.125) {freed, not reclaimed};
  \draw[gap]  (8.6,0) rectangle (9.0,0.25); \node[right] at (9.05,0.125) {wasted by a wrap};
  \draw[fill=gray!6] (12.0,0) rectangle (12.4,0.25); \node[right] at (12.45,0.125) {free};
\end{scope}
\end{tikzpicture}
\caption{A ring allocator in six snapshots, from left to right. The pool
  is one contiguous region, addresses grow upwards, and allocations are
  any size, packed back to back. New allocations go at the \emph{head}
  (orange arrow, right), and space is reclaimed only at the \emph{tail}
  (blue arrow, left), the start of the oldest allocation still live.
  Freeing out of order is allowed but reclaims nothing yet: in (c), B is
  freed before A, so its space stays occupied until A is freed too in (d),
  and the tail then moves past both. When a request does not fit above the
  head, the head wraps to the bottom of the pool, and the rest of the top
  is wasted until the tail passes it (e). \texttt{deallocate\_all} resets
  both cursors at once (f).}
\label{fig:ring-alloc}
\end{figure}
